\section{Implementation}
\label{sec:carp:impl}


In this section we describe CARP's implementation and its relationship to scalability, memory footprint, and runtime overhead.

CARP \cite{Jain2024:CARP} consists of \mytexttilde10,000 lines of C/C++ code. Application writes are intercepted using a dynamically preloaded shared library. The CARP API can also be called directly without using a preload library. CARP uses two instances of the DeltaFS shuffle service~\cite{Zheng2017} for communications: the \emph{data} instance batches messages into 32KB buffers for high throughput, while the \emph{control} instance is used for renegotiation control messages and uses no batching for low latency. The shuffle service is built using the Mochi \cite{Ross2020:Mochi} Mercury RPC library~\cite{Souma2013:MercuryRPC}. We use RPCs as both shuffling and renegotiation are asynchronous operations better suited to RPC semantics (it also allows us to avoid complications with multiple concurrent MPI communicators). CARP creates its own threads for network communication, RPC callbacks, and asynchronous compaction in KoiDB. This setup allows straggler ranks to use multiple cores and provides some tolerance for load imbalance.

\textbf{Scalability.}
CARP is designed to scale by composing scalable operations --- the shuffle service has been demonstrated to scale to 131,072 ranks~\cite{Zheng2018:DeltaFS}, and the renegotiation protocol is designed as a reduction operation with a logarithmic scaling factor (evaluated in \S\ref{sec:carp:eval_rtpscal}). Each instance of CARP's storage backend (KoiDB) is an independent instance and can scale trivially. By default, CARP creates one file per rank, but at larger scales, the total number of files can be reduced (if needed) by having a subset of ranks be shuffle receivers.

We call CARP's scalable implementation of renegotiation the \emph{Tree-based Renegotiation Protocol (TRP)}. A naive implementation of renegotiation, as described in \S\ref{sec:carp:des_reneg}, would require directly collecting all ranks' pivots on one rank before computing the new partition table. Such an implementation will have limited scalability, as its memory and network communication footprint will scale linearly with ranks. Our TRP implementation, on the other hand, uses a reduction tree-based design~\cite{Raben2004:OptimizationCollectiveReduction, Raben2004:MoreEfficientReduction} which scales logarithmically.

TRP is built on the observation that pivot unions (\S\ref{sec:carp:des_hist}), being associative and commutative have all the properties of a lossy reduction operation. Pivots represent distributions and can be merged in any order, but the pivots lose some information in the process. To limit the impact of this lossiness, TRP employs a shallow tree hierarchy with a large fan-out (depth of 3, fan-out of up to 64). The tree's leaf layer consists of all CARP ranks, while intermediate layers consist of equally spaced ranks. Subsets of pivots are reduced on intermediate layers, and a final reduction happens at the root.

\textbf{Memory Footprint.}
As CARP runs in application processes, it needs to have a low memory footprint to avoid competing with the application for memory.  CARP achieves this by streaming application data directly to the shuffle pipeline.  The primary source of memory overhead is due to buffers used to batch I/O for network and storage efficiency.  To illustrate this overhead, we use a run with 4096 ranks and default CARP parameters that results in 27MB per rank\footnote{Each rank uses 2MB for shuffle RPC buffers~\cite{Zheng2018:DeltaFS}, 24MB for two KoiDB memtables, 16KB for 4096*4B partition table entries, 16KB for partition shuffle counts, and 32KB for 512-entry OOB buffer with 64B records.}. This is less than 1\% of the per-core memory budget on the LANL Trinity supercomputer~\cite{Vigil2014:Trinity}.

\textbf{Runtime Overhead.}
CARP only runs during application I/O, so it does not impact the compute phase.  Data shuffling and triggered renegotiations are potential sources of additional runtime overhead.  Shuffle bandwidth increases with the number of writes and quickly outstrips storage bandwidth in our experiments. Thus, the shuffle service itself does not have any runtime overhead.  CARP could underutilize storage by pausing I/O for renegotiation, but our experiments show that shuffle receivers buffer enough outstanding writes to keep storage busy and mask this overhead.  If needed, CARP can continue routing data using the old partition table while a renegotiation is underway, however we have not found this necessary in our experiments.

\begin{figure}[t]
  \centering
  \begin{subfigure}[b]{0.48\textwidth}
  \figktqlat
  \end{subfigure}
  \hspace{1em}
  \begin{subfigure}[b]{0.48\textwidth}
  \figktrt
  \end{subfigure}
  \capfigkt
\end{figure}
